\documentclass[UTF8,a4paper,10pt]{ctexart}
\usepackage{ipara-handbook}
\usepackage{booktabs,tabularx}
\begin{document}
\section*{H-B04\quad 连续无限累积与离散无限累积}
\textbf{边界：}本桥翻译共同的“有限截断到无限对象”构造，不把积分判别法和级数判别法混成一套公式。
\begin{tcolorbox}[title=共同骨架]
有限截断 $\longrightarrow$ 尾部稳定（Cauchy）$\longrightarrow$ 比较/抵消
$\longrightarrow$ 取极限 $\longrightarrow$ 记录收敛类型。
\end{tcolorbox}
\begin{tabularx}{\linewidth}{@{}lXX@{}}
\toprule
对象 & 截断 & 尾部检查\\ \midrule
反常积分 & $\int_a^R f,\ R\to\infty$ 或 $\int_{a+\varepsilon}^b f,\ \varepsilon\to0^+$ & 任意两个截断的差趋零；奇点与无穷远分开\\
数项级数 & $S_N=\sum_{n=1}^{N}a_n,\ N\to\infty$ & $\sum_{n=m}^{N}a_n$ 对 $m,N$ 足够大趋零\\
\bottomrule
\end{tabularx}
\subsection*{1.\ 同一个尾部语言，不同的原子单位}
级数研究离散尾段
\[
\sum_{k=m}^{n}a_k,
\]
反常积分研究连续尾段
\[
\int_u^v f(x)\,dx.
\]
两者的 Cauchy 语言相同：任意足够远/足够靠近奇点的尾段贡献都必须统一变小。但二者的“原子单位”不同，因此不能把必要条件机械互译。

数项级数收敛必有
\[
a_n\to0,
\]
因为单项就是相邻两个部分和之差。反常积分收敛却\textbf{不要求} $f(x)\to0$。例如可以在每个整数 $n$ 附近放一个高度为 $1$、底宽为 $2^{-n}$ 的连续三角尖峰，并令其他地方为 $0$。总面积至多
\[
\sum_{n=1}^{\infty}2^{-n}<\infty,
\]
所以积分可收敛，但 $f(n)=1$ 对所有 $n$ 成立，故 $f(x)$ 不趋于 $0$。

\begin{tcolorbox}[title=核心 Anti-Bridge]
\[
\boxed{\sum a_n\text{ 收敛}\Rightarrow a_n\to0,
\qquad
\int_A^\infty f\text{ 收敛}\not\Rightarrow f(x)\to0.}
\]
共同的是“\textbf{尾部总贡献}趋零”，不是“局部高度/单项一定趋零”。只有再加单调、连续等结构时，积分收敛才可能强迫更强的点态结论。
\end{tcolorbox}

\subsection*{2.\ 积分判别：把离散矩形和连续面积夹在一起}
若 $f$ 在 $[1,\infty)$ 非负、连续且单调递减，则
\[
\sum_{n=1}^{\infty}f(n)\ \text{与}\ \int_1^\infty f(x)\,dx
\]
同敛散，并可用矩形面积得到尾项夹逼：
\[
\int_{N+1}^{\infty}f(x)\,dx
\le
\sum_{n=N+1}^{\infty}f(n)
\le
\int_N^{\infty}f(x)\,dx.
\]
这条 Bridge 的本质是“单调性让离散采样矩形稳定地包住连续面积”。条件缺一时，不能只因 $f(n)$ 与 $f(x)$ 长得相同就宣告同敛散；应回到各自 Topic 的比较或 Cauchy 判别。

\subsection*{3.\ 绝对与条件}
\[
\sum \lvert a_n\rvert<\infty\Rightarrow\sum a_n\text{ 收敛},
\]
但反向不成立；积分同样要区分绝对收敛与条件抵消。绝对收敛允许更多安全的拆分、重排/换序，而条件收敛依赖抵消结构。

这里必须特别阻断“对称抵消 = 普通收敛”。例如
\[
\operatorname{PV}\int_{-1}^{1}\frac{dx}{x}=0
\]
只表示对称截断的 Cauchy 主值存在；普通反常积分要求 $0$ 左右两段分别收敛，而它们都发散。离散世界中条件收敛级数的重排可以改变和，连续世界中不同截断路径/主值也可能改变结果；两者共享的是“抵消对操作顺序敏感”这一结构，不是共享同一个定义。

\subsection*{4.\ 连续分部积分与离散 Abel 变换：同一抵消机制的两种语言}
积分中的分部积分
\[
\int u\,dv=uv-\int v\,du
\]
与级数中的 Abel 变换
\[
\sum_{n=1}^{N}a_nb_n
=A_Nb_N+\sum_{n=1}^{N-1}A_n(b_n-b_{n+1}),
\qquad A_n=\sum_{k=1}^{n}a_k,
\]
都把“振荡对象 $\times$ 缓慢变化权重”拆成\textbf{边界项 + 累积对象乘权重变化}。因此 Dirichlet 型判敛在连续/离散两边都有同一母机制：
\begin{itemize}
  \item 振荡一侧的原函数/部分和保持有界；
  \item 权重单调趋零；
  \item 边界项消失，剩余变化量可控制。
\end{itemize}
这条 Bridge 只解释机制翻译；具体积分 Dirichlet/Abel 条件归 Topic05，级数版本归 Topic10。

\subsection*{5.\ 最小例子与调用}
取 $f(x)=x^{-p}$。当 $p>0$ 时它在 $[1,\infty)$ 连续、非负且递减，故
\[
\sum_{n=1}^{\infty}\frac1{n^p}\quad\text{与}\quad
\int_1^\infty x^{-p}\,dx
\]
同敛散，临界点都是 $p=1$。题面在反常积分与正项级数之间比较尾部时，先核对积分判别资格。

\subsection*{6.\ 迁移时的反桥}
连续变量的奇点位置、离散序列的尾部索引、单调性和非负性分别核对。共享“尾部稳定”思想，不意味着可把某个比较常数、必要条件或端点估计直接复制。尤其牢记：
\begin{itemize}
  \item 级数的单项趋零必要条件不能直接搬成反常积分的 $f(x)\to0$；
  \item 积分判别需要非负、连续、单调等资格，不是“和与积分总能互换”；
  \item 条件抵消下，级数重排与积分主值都是额外操作，不能冒充原定义；
  \item 比较的是尾部总贡献；局部峰值、单项和采样点只是不同表示中的局部信息。
\end{itemize}
\begin{tcolorbox}[title=检查句]
先写截断对象和极限方向，再写尾部标准；最后才选择比较、积分判别或绝对收敛工具。
\end{tcolorbox}
\end{document}
